\documentclass[11pt, a4paper]{article}
\input{Config.tex}
\usepackage{hyperref}
\renewcommand{\thesection}{}
\renewcommand{\thesubsection}{}
\renewcommand{\thesubsubsection}{}





\title{Problemes de the Rising Sea}
\author{Bernat Esteve Sagarra}

\begin{document}
\pagestyle{fancy}
% 
\lhead{Bernat Esteve}
\rhead{Problemes de The Rising Sea}
\section*{Capítol 1}
\begin{tcolorbox}[title=\subsection{Exercici 1.1.B}]
    
    If $\mathcal{A}$ is an object in a category $\CCat$, show that the invertible elements of $\Mor(\mathcal A, \mathcal A)$ form a group (called the automorphism group of $\mcal A$, denoted $\Aut(\mcal A)$).\\
    What are the automorphism groups of the objects in Examples 1.1.2 and 1.1.3? Show that two isomorphic objects have isomorphic automorphism groups. (For readers with a topological background: if $X$ is a topological space, then the fundamental groupoid is the category where the objects are points of $X$, and the morphisms $x \to y$ are paths from $x$ to $y$, up to homotopy. Then the automorphism group of $x_0$ is the (pointed) fundamental group $\pi_1(X, x_0)$. In the case where $X$ is connected, and $\pi_1(X)$ is not abelian, this illustrates the fact that for a connected groupoid — whose definition you can guess — the automorphism groups of the objects are all isomorphic, but not canonically isomorphic.)
\end{tcolorbox}
\subsubsection{Exercici 1.1.B}
Els morfismes invertibles en la categoria dels $\SSet$ són les bijeccions que 
van d'un conjunt a ell mateix; i en la categoria dels $\VVec_\KK$, els automorfismes, són endomorfismes amb determinant no nul.
\end{document}
% ---
% \newpage
% \begin{tcolorbox}[title=\subsection{Exercici XXX}]
% \end{tcolorbox}
% \subsubsection{Exercici XXX}
